% TOP_PROBLEM: {"id": 1269, "title": "Waring's Problem: Exact Values", "category_id": 1, "category": {"id": 1, "name": "number_theory", "display_name": "Number Theory", "description": "Properties of integers, prime numbers, Diophantine equations.", "slug": "number-theory", "order_index": 1, "created_at": "2026-07-31T15:26:25.670Z"}, "status": "open"}
% FIRST_POSTED: null
% ENTRY_KIND: statement_only
% Imported from: batch05.tex; UnsolvedMath ID 1269
\documentclass[11pt]{article}
\usepackage[margin=25mm]{geometry}
\usepackage{fontspec}
\setmainfont{Latin Modern Roman}
\usepackage{amsmath,amssymb,mathtools}
\usepackage{unicode-math}
\setmathfont{Latin Modern Math}
\usepackage{xurl,hyperref}
\hypersetup{colorlinks=true,urlcolor=blue}
\setcounter{secnumdepth}{0}
\setlength{\parindent}{0pt}
\setlength{\parskip}{5pt}
\setlength{\emergencystretch}{3em}
\newcommand{\sref}[2]{\hyperlink{source:#1:#2}{[S#2]}}
\newcommand{\cataloguescope}[1]{\paragraph{Recorded target.}#1}
\begin{document}
% UnsolvedMath ID 1269; source catalogue code NT-062
\setcounter{section}{1268}
\section{Waring's Problem: Exact Values}\label{1269}
{\small\sffamily UnsolvedMath ID 1269 · Number Theory}
\subsection{Definitions and mathematical statement}

\paragraph{Shared notation.}$\mathbb N=\{1,2,\ldots\}$, and $\mathbb Z,\mathbb Q,\mathbb R,\mathbb C$ denote the integers, rationals, reals and complex numbers. Any other conventions are specified locally.


Let $k\in\mathbb N$, and let $\mathbb Z_{\ge0}=\{0,1,2,\ldots\}$. Define $R_{k,s}(n)$ to mean that there exist $x_1,\ldots,x_s\in\mathbb Z_{\ge0}$ with $n=x_1^k+\cdots+x_s^k$. Zero summands permit representations using at most $s$ positive powers. Set
\[g(k)=\min\{s\in\mathbb N:\forall n\in\mathbb N, R_{k,s}(n)\},\]
\[G(k)=\min\{s\in\mathbb N:\exists N\in\mathbb N\ \forall n\ge N, R_{k,s}(n)\}.\]
Thus $g$ requires every positive integer, while $G$ allows finitely many exceptions. Hilbert's theorem ensures these minima exist for $k\ge2$, and $g(1)=G(1)=1$.
\paragraph{Statement recorded in the supplied source.}
Determine the exact integers $g(k)$ and $G(k)$ for every integer exponent $k\ge2$ (with the trivial $k=1$ values as above). Both the universal and sufficiently-large representation functions are included. \sref{1269}{2} \sref{1269}{3}
The two functions have different quantifiers: determining a formula for $g(k)$ alone does not determine $G(k)$.
\subsection{Short English statement}
Determine exactly how many nonnegative integral kth powers suffice for every positive integer and how many suffice after excluding finitely many small integers, for each exponent.
\subsection{Sources}
\begin{description}
\item[\hypertarget{source:1269:1}{S1}] ULAM.AI and the UnsolvedMath Contributors, UnsolvedMath: A Curated Collection of Open Mathematics Problems, record 1269 (NT-062). CC BY 4.0. \url{https://huggingface.co/datasets/ulamai/UnsolvedMath}.
\item[\hypertarget{source:1269:2}{S2}] Michel Waldschmidt, \emph{Open Diophantine Problems}, arXiv:math/0312440 (version dated January 24, 2004), Section 2.2, pp.\ 16--17, discussion of Waring\textquotesingle s problem. \url{https://arxiv.org/pdf/math/0312440}.
\item[\hypertarget{source:1269:3}{S3}] J\"org Br\"udern and Trevor D. Wooley, \emph{On Waring\textquotesingle s problem for larger powers}, arXiv:2211.10380v1 (2022), Section 1, p.\ 1. \url{https://arxiv.org/pdf/2211.10380v1}.
\end{description}
\label{end:1269}
\end{document}
